\section{Methodology}
\label{sec:method}

\Cref{fig:pipeline} shows the pipeline. It has two separate views of the same
users (\cref{sec:method-views}), per-node membership vectors and disagreement scores
(\cref{sec:method-scores}), a detection check (\cref{sec:method-detect}) and a spread
simulation that tests interventions (\cref{sec:method-block}).

\begin{figure}[t]
\centering
\begin{tikzpicture}[
  box/.style={draw, rounded corners, align=center, minimum height=1.0cm,
              text width=3.0cm, font=\small},
  arr/.style={-{Latex[length=2mm]}, thick},
  node distance=0.7cm and 0.9cm]
  \node[box] (graph) {Social graph $A$\\(symmetrised, LCC)};
  \node[box, below=of graph] (text) {User text $x_i$\\(+ \texttt{has\_metadata})};

  \node[box, right=of graph] (leiden) {Topology view\\Leiden $\times R$ $\to t$};
  \node[box, right=of text] (clf) {Content view\\cross-fitted text\\classifier $\to p_C$};

  \node[box, right=of leiden] (spread) {Label spreading\\$\to p_T$};
  \node[box, below=of spread] (scores) {Scores: misfit,\\disagreement, baselines};

  \node[box, below=1.0cm of clf] (detect) {Detection check\\synthetic + TIMME};
  \node[box, right=of detect] (block) {Spread simulation\\node / link blocking,\\regime map};

  \draw[arr] (graph) -- (leiden);
  \draw[arr] (text) -- (clf);
  \draw[arr] (leiden) -- (spread);
  \draw[arr] (leiden) -- node[left, font=\scriptsize] {labels $t$} (clf);
  \draw[arr] (spread) -- (scores);
  \draw[arr] (clf) -- (scores);
  \draw[arr] (scores) -- (block);
  \draw[arr] (scores.south west) -- (detect.north east);
\end{tikzpicture}
\caption{Pipeline. The topology view uses edges only. The content view receives no
edges or graph-derived features as input; the topology labels $t$ are used only as
its training targets, with cross-fitting so that no node's own label informs its
prediction.}
\label{fig:pipeline}
\end{figure}


\subsection{Problem formulation}
\label{sec:method-problem}
Let $G=(V,E)$ be a social graph with $n=|V|$ users and adjacency matrix $A$. Every
user $i$ has a text vector $x_i$ (\cref{sec:setup-features}), which is zero if the
user has no text. We treat ties as undirected (requirement R5,
\cref{sec:setup-requirements}). The aim is to:
\begin{enumerate}
  \item assign every user a structural community $t_i \in \{1,\dots,K\}$ from $A$ alone;
  \item estimate two soft memberships over those $K$ communities, a structural one
        $p_T(i)$ and a content one $p_C(i)$, where $p_C(i)$ depends only on text;
  \item score how far apart they are; and
  \item decide whether blocking high-scoring nodes or links reduces spread from an
        \emph{unknown} source more than structural targeting does, for a fixed budget.
\end{enumerate}

\subsection{Design evolution}
\label{sec:method-evolution}
The first design followed AAS closely. Each node would be embedded from text and
topology together with a self-supervised GNN (BGRL, DGI or MVGRL on a shared two-layer
GCN encoder), and the embedding would be clustered by a density method (HDBSCAN) and a
first-neighbour method (FINCH). Ambiguous regions would come from where those two
clusterings partly overlap. Diagnostics during development showed that the clusters of
these embeddings mainly sorted nodes by degree rather than by community
(\cref{tab:gnn-diag}). On TIMME, for all three objectives, the two-way $k$-means split
of the embedding and its first principal component were both strongly correlated with
log-degree, while the six-way $k$-means partition had negative modularity on the
social graph and almost no mutual information with party. Party information was not
missing from the embeddings: a linear probe predicted party well above the majority
baseline. It was present, but the cluster structure followed degree, not communities.
These numbers are a development diagnostic of the discarded design, shown only because
they motivated the redesign. There was also a design flaw: a content view derived from
an embedding that has already seen the graph cannot disagree with the graph in a
meaningful way. The GNN stage was therefore replaced by the two separate views below. The AAS principle of looking where
complementary views disagree is kept.

\begin{table}[t]
\centering
\caption{Development diagnostic of the discarded GNN design on TIMME (20,811 nodes,
1,206 with party labels). Clusters are $k$-means on L2-normalised embeddings. Shown
only to justify the redesign (\cref{sec:method-evolution}); these are not results of
the current pipeline.}
\label{tab:gnn-diag}
\small
\begin{tabular}{@{}l r r r@{}}
\toprule
Diagnostic & BGRL & DGI & MVGRL \\
\midrule
$|r|$, $k=2$ split vs.\ log-degree           & 0.58   & 0.60   & 0.67   \\
$|r|$, first principal component vs.\ log-degree & 0.66 & 0.70  & 0.78   \\
Modularity of the $k=6$ partition            & $-0.050$ & $-0.152$ & $-0.104$ \\
Party NMI of the $k=6$ partition             & 0.060  & 0.018  & 0.002  \\
Linear party probe accuracy (5-fold)$^{*}$   & 0.943  & 0.934  & 0.780  \\
\bottomrule
\end{tabular}

\smallskip
\footnotesize
$^{*}$Logistic regression on the embedding, labelled users only; always predicting
the majority party gives 0.534.
\end{table}


\subsection{Preprocessing}
\label{sec:method-pre}
For every dataset we remove self-loops, make every edge two-way and keep the
\emph{largest connected component} (a node with no path to the rest of the network
cannot pass a message on). Users without text are \emph{kept}, because removing them
would also remove their edges, and some of those edges may be bridges. Such users get
a zero text vector and a binary \texttt{has\_metadata} flag.

\subsection{Two separate views}
\label{sec:method-views}

\paragraph{Topology view.}
Leiden~\citep{traag2019leiden} with the modularity objective is run on $G$ with $R$
different random seeds. The run with the highest modularity is the reference partition
$t$. Communities containing fewer than 1\% of the nodes are set aside ($t_i=-1$): they
are never used as training targets or as seeds for spreading. The remaining $K$
communities form the label space. All $R$ runs are kept for the consensus baseline.

\paragraph{Content view.}
Content membership is estimated with \emph{cross-fitting}. The users with $t_i \ge 0$
are divided into $F$ stratified folds. For each fold, a multinomial logistic
regression (with standardised inputs and $L_2$ regularisation) is trained on the other
folds to predict $t$ from the text vector $x$, and is then applied to the held-out
fold. The result is
\begin{equation}
  p_C(i) = \Pr\bigl(t_i = \cdot \mid x_i\bigr) \in \Delta^{K-1},
\end{equation}
where every $p_C(i)$ comes from a model that never saw user $i$'s own label. Users with
$t_i=-1$ are scored by every fold model and the predictions are averaged. The
classifier receives no edges or graph-derived features as input; the topology labels
$t$ are used only as training targets, with cross-fitting so that no node's own label
informs its prediction. The two views are therefore separate in their inputs, but not
fully independent: the content view learns which text goes with which graph
community.

\paragraph{Structural soft membership.}
The hard partition $t$ is turned into soft memberships by label
spreading~\citep{zhou2004learning}. With $Y$ the one-hot matrix of $t$ (zero rows for
$t_i=-1$) and $S = D^{-1/2} A D^{-1/2}$,
\begin{equation}
  F^{(m+1)} = \alpha\, S F^{(m)} + (1-\alpha)\, Y, \qquad F^{(0)} = Y,
\end{equation}
iterated a fixed number of times. Each row is then normalised to give $p_T(i)$. Users
deep inside a community end up with an almost one-hot $p_T(i)$; users on a boundary
end up with mixed rows.

\subsection{Per-node scores}
\label{sec:method-scores}
Three scores use the two views:
\begin{align}
  \text{misfit}_i &= 1 - p_C(i)_{t_i}, \\
  \text{disagreement}_i &= \mathrm{JS}_2\bigl(p_T(i)\,\Vert\,p_C(i)\bigr), \\
  \text{ambiguity}^T_i &= H\bigl(p_T(i)\bigr) / \log K.
\end{align}
Misfit measures how little the text supports the node's own community.
Disagreement is the base-2 Jensen--Shannon divergence~\citep{lin1991divergence}
between the two memberships, which lies in $[0,1]$. Structural ambiguity is the
normalised entropy of the structural membership and uses no text.
For $t_i=-1$, misfit is defined as $1-\max_k p_C(i)_k$.

\paragraph{Baselines.}
These are the scores the proposed ones have to beat. Each captures one alternative
explanation of what disagreement might be measuring:
\begin{itemize}
  \item \textbf{Consensus entropy} (structure only): every Leiden run is aligned to $t$
        with the Hungarian algorithm, and the score is the normalised entropy of each
        node's label across runs~\citep{riolo2020consistency}.
  \item \textbf{Participation coefficient}~\citep{guimera2005functional}:
        $P_i = 1-\sum_c (h_{ic}/d_i)^2$, where $h_{ic}$ is the number of $i$'s
        neighbours in community $c$ and $d_i$ is $i$'s degree.
  \item \textbf{Content outlier}, following \citet{gao2010outliers}:
        $1-\cos(x_i, \bar{x}_{t_i})$, the cosine distance from the centroid of the
        node's community's text.
  \item \textbf{Degree} and \textbf{PageRank}, to check that a score is not simply
        degree in disguise.
\end{itemize}

\subsection{Checking that the scores detect real conflicts}
\label{sec:method-detect}

\paragraph{Synthetic benchmark with planted conflicts.}
Synthetic networks have a known answer. The generator is a degree-corrected block
model in the style of LFR~\citep{lancichinetti2008benchmark,karrer2011stochastic}.
Degrees follow a power law. A share $\mu$ of each node's edge stubs is matched across
communities (the mixing parameter), and the rest is matched inside the node's own
community. Each node's text is its community's centroid plus Gaussian noise of scale
$\sigma$, normalised to unit length. A small share of nodes is \emph{planted}: they
get the text centroid of a different community, and these are the true conflicts.
Real text also varies for reasons that have nothing to do with community, so an
optional \emph{nuisance} block adds random topics and a writing-style direction
proportional to log-degree. The full pipeline (Leiden, cross-fitted classifier,
scores) is run on each network. Every score is evaluated by ROC AUC for recovering
the planted nodes, both overall and within degree terciles (so a score cannot win
just by tracking degree).

\paragraph{Real-data check on TIMME.}
TIMME has party labels for part of the users (\cref{sec:setup-data}). A labelled user
counts as a \emph{party misfit} when their party differs from the majority party of
their topology community. Only communities with at least 20 labelled users and at
least 80\% one party are used, because a mixed community has no ``wrong'' side. Each
score is evaluated by AUC for predicting party misfits, overall and within degree
terciles, with a one-sided Mann--Whitney test. As a null, the text rows are shuffled
across users and all scores are recomputed. A real score has to beat its own
shuffled-text version.

\subsection{Testing interventions with spread simulation}
\label{sec:method-block}

\paragraph{Spread model.}
We use live-edge Independent Cascade~\citep{kempe2003maximizing}. An edge $u\to v$
transmits with probability $p_{uv}$ under one of two schemes:
\begin{itemize}
  \item \emph{weighted cascade}, $p_{uv}=1/\deg(v)$;
  \item \emph{uniform}, $p_{uv}=c\cdot p_c$, where
        $p_c = \langle k\rangle / (\langle k^2\rangle - \langle k\rangle)$ is the
        epidemic threshold of a network with this degree
        sequence~\citep{newman2002spread}, and $c$ is swept over several multiples.
\end{itemize}

\paragraph{Unknown source and common random numbers.}
In every simulation the seeds are drawn uniformly from nodes that \emph{no} strategy
blocks at the largest budget, so no strategy can win by deleting the source. Every
strategy and budget sees the same live-edge samples and the same seeds (common random
numbers), so differences between strategies are not Monte Carlo noise.

\paragraph{Outcomes.}
For each run we record the \emph{infected share} (seeds excluded) and the
\emph{cross-community share}, i.e.\ infected users outside every seed's community.
Each strategy is reported as the percentage of spread it removes compared with no
blocking.

\paragraph{Content-aware spread.}
Under selective exposure~\citep{delvicario2016spreading} a message is accepted more
readily by users whose content matches the community it came from. Each cascade
starts in one community $A$ (seeds drawn inside $A$), and every edge probability is
reweighted:
\begin{equation}
  p'_{uv} = \min\bigl(1,\; p_{uv}\cdot K\cdot a_v(A)\bigr),
\end{equation}
where $a_v \in \Delta^{K-1}$ is node $v$'s affinity to each community's content. The
multiplier $K\,a_v(A)$ averages approximately 1 over receivers, so the overall
transmission rate stays roughly the same. This is only approximate: the affinities
average exactly $1/K$ across receivers only when the communities are of equal size,
and the $\min(1,\cdot)$ clipping lowers the effective rate further. On
TIMME, $a_v$ comes from a separate cross-fitted classifier trained on a \emph{different}
text source (GloVe tweet averages) from the one used for scoring (MiniLM). This keeps
the spread model from reusing the scorer's own predictions.

\paragraph{Node-blocking strategies.}
Given a budget $k$, each strategy removes its top-$k$ nodes. Structural strategies:
random, degree, adaptive degree (recomputed after each removal), PageRank, sampled
betweenness, Collective Influence at radius 1~\citep{morone2015influence},
participation coefficient, and participation, consensus entropy and structural
ambiguity each multiplied by PageRank. Content-based strategies: misfit,
disagreement, content outlier (alone and multiplied by PageRank), and a
\emph{gateway} score
\begin{equation}
  g_i = \sum_{A} h_{iA}\; a_i(A)\; \bigl(d_i - h_{iA}\bigr),
\end{equation}
which estimates how much of community $A$'s content node $i$ can carry out of $A$.
In words: for each community $A$, the number of $i$'s neighbours in $A$, weighted by
how well $i$'s content matches $A$, is multiplied by the number of $i$'s neighbours
outside $A$, and these products are summed over communities. The
same formula with $a_i(A)=1/K$ (``structural gateway'') removes the text, so comparing
the two isolates what the content adds.

\paragraph{Link-level intervention.}
Instead of removing accounts, we can add friction to individual links (both directions
of an undirected link are cut together). Link scores: random, degree product, PageRank
product, a structural bridge score (cross-community links ranked by
$\mathrm{PR}_u\mathrm{PR}_v$), and two content-aware variants. The first ranks
cross-community links by
$\max\bigl(\mathrm{PR}_u\, p_C(v)_{t_u},\ \mathrm{PR}_v\, p_C(u)_{t_v}\bigr)$, i.e.\ an
influential sender times a receptive receiver. The second multiplies the structural
bridge score by that receptivity.

\paragraph{Regime map.}
Real networks fix the community strength and the size of the content effect, so a
single dataset cannot show when content-aware targeting helps. We therefore sweep
these factors on synthetic networks:
\begin{itemize}
  \item the mixing $\mu$;
  \item the agreement effect $\beta$, which sets the affinity
        $a_v(A) \propto \exp\bigl(\beta\,[c_v = A]\bigr)$ for node $v$'s \emph{true}
        content community $c_v$ ($\beta=0$ is content-blind spread);
  \item the share of mismatched nodes; and
  \item where the mismatched nodes are placed: uniformly at random, or in proportion
        to degree (``hub'').
\end{itemize}
Scores are computed from \emph{observed}, noisy text through the full pipeline, while
spread uses the true content. For each setting we report the paired gain, in
cross-community spread removed, of each content-aware strategy over the better of
PageRank and betweenness, with a 95\% confidence interval.

\subsection{Planned extension: active learning}
\label{sec:method-al}
AAS uses disagreement to choose must-link and cannot-link pair queries for an oracle
and then refines the clustering under those constraints. A version of that loop, with
queries drawn from high-disagreement nodes and a simulated oracle built from TIMME
party labels, is planned as an extension and depends on discussion with the
supervisor. It is not part of the work described in this report.

\subsection{Limitations}
\label{sec:method-limits}
Three limitations of the design are known in advance.
\begin{itemize}
  \item \textbf{Content-aware spread may favour content scores by construction.} The
        receiver affinity $a_v$ is itself a text-to-community classifier trained on
        the same topology labels $t$. Strategies that use $p_C$ may therefore look
        better under content-aware spread partly because the spread model and the
        scores were built the same way. Using GloVe for the spread model and MiniLM
        for the scores reduces this overlap, but both come from the same users'
        tweets, so it does not remove it.
  \item \textbf{Sampling of the TIMME crawl.} TIMME was crawled around seed accounts
        (politicians), so some apparent bridges may be artefacts of how the data was
        collected rather than real crossing points (\cref{sec:setup-data}).
  \item \textbf{Simulated spread on a static graph.} All interventions are evaluated
        by simulation on a static, historical snapshot of each network, with assumed
        Independent Cascade probabilities rather than probabilities estimated from
        observed cascades. This limits how strongly any finding can be claimed to hold
        for real-world spread.
\end{itemize}
